\chapter{Experimental Results}
\label{ch:results}

% One subsection tree per dataset. Each dataset section follows the same
% internal shape for comparability: summary table -> per-model discussion ->
% dataset-specific takeaway. Cross-dataset synthesis (which paradigm wins
% where, and why) is saved for the Discussion/Conclusion chapter, not here --
% keep this chapter to "what happened", not "what it means for the thesis".

\section{Credit Card Fraud Detection (Dataset A)}
% Source docs: data/results/results_credit_card/*_report.txt

\subsection{Overall Comparison}
% Summary table: model x {ROC-AUC, Average Precision}, all 8 models
% (IForest, OC-SVM, MLP AE, Predictive LSTM, LSTM AE, Bottleneck Transformer,
% Causal Transformer... note: no causal transformer implementation exists
% for credit card, per architecture_diffs.md).

\subsection{Discussion}
% Expected narrative: classical / reconstruction models >= sequential models,
% since the 29 PCA features have no meaningful temporal order. Note the
% "pseudo-sequence" caveat for LSTM/LSTM-AE/Transformer on this dataset.

\section{Yahoo Webscope S5 (Dataset B)}
% Source docs: yahoo/results_notes.md, data/results/results_yahoo/*_report.txt

\subsection{A1 -- Point Anomalies (Real Server Metrics)}
\subsection{A3 -- Contextual/Seasonal Anomalies}
\subsection{A4 -- Level-Shift Anomalies}
% Per-benchmark: summary table, then the key finding already established in
% results_notes.md (e.g. A3: Predictive LSTM dominates via implicit phase
% tracking; A4: expected reconstruction-model reversal did not materialize).

\subsection{Trivial Baseline Sanity Check}
% \|w\|_2 baseline (Kim et al. 2022, Case 2) confirming A3 gains are real
% signal, not window-magnitude artifacts.

\subsection{Discussion}
% The 2x2 paradigm x architecture finding: predictive necessary but not
% sufficient; LSTM's recurrent hidden state as phase accumulator vs. causal
% attention. Note A2 as future/incomplete work if it stays that way.

\section{CIC-IDS2017 (Dataset C)}
% Source docs: cicids2017/results_analysis.md

\subsection{Flat vs. Temporal vs. Host-Context Features}
% Summary table across all 13 models. Lead with the context-feature finding:
% PortScan AP 0.07-0.28 (flat/temporal) -> 0.84 (AE+Context).

\subsection{Per-Attack-Type Breakdown}
% Per-class AP heatmap discussion (DDoS, PortScan, slow DoS, etc.) --
% which model/feature-set handles which attack family.

\subsection{Discussion}
% Reconstruction-based > prediction-based (attack traffic more predictable
% than benign); context dominates; temporal windowing adds limited value
% given arbitrary intra-window flow ordering. Flag known re-run needed:
% Bottleneck Transformer with corrected LATENT_DIM=32.
